\chapter{Octalysis}
\label{ch:octa}
Punteggio rilevato dal radar \texttt{docs/assets/Octa\_Analysis.jpg}: \textbf{169}.

\section{Gli 8 Core Drive}
\begin{table}[h]
\centering
\begin{tabular}{@{}clr@{}}
\toprule
CD & Drive & Slay \\
\midrule
CD1 & Epic Meaning \& Calling & 5 \\
CD2 & Development \& Accomplishment & 5 \\
CD3 & Empowerment of Creativity \& Feedback & 7 \\
CD4 & Ownership \& Possession & 3 \\
CD5 & Social Influence \& Relatedness & 2 \\
CD6 & Scarcity \& Impatience & 2 \\
CD7 & Unpredictability \& Curiosity & 7 \\
CD8 & Loss \& Avoidance & 2 \\
\midrule
\multicolumn{2}{@{}l}{Totale} & 33 ($\rightarrow$ score 169 normalizzato) \\
\bottomrule
\end{tabular}
\caption{Core Drive Octalysis di Slay the Roadmap.}
\label{tab:octa}
\end{table}

Il profilo è sbilanciato su CD3 e CD7 (creatività del deck e curiosità dei
quiz/boss) con CD1/CD2 medi (missione roadmap, XP/badge) e CD4--CD6/CD8
bassi: il gioco premia l'esplorazione più che il possesso o la pressione
sociale.

\section{Mapping US $\rightarrow$ Core Drive}
\begin{itemize}[nosep]
\item US-01 (roadmap ad albero): CD1 (chiamata della missione), CD2 (progresso visibile).
\item US-02 (quiz 80\%): CD2 (mastery), CD3 (feedback immediato).
\item US-03 (ricompense/deck): CD3 (combinazioni creative), CD4 (possesso carte).
\item US-04 (boss fight): CD7 (quiz adattivi imprevedibili), CD8 (rischio sconfitta).
\item US-05 (persistenza): CD4 (proprietà del progresso), CD2 (continuità).
\end{itemize}
CD5 e CD6 restano scoperti: le leaderboard sociali e la scarsità (energia/vite)
sono mock statici (3/5 vite, 5/5 energia) e costituiscono un lavoro futuro.
